\chapter*{Abstract}
\addcontentsline{toc}{chapter}{Abstract}

Online education produces hundreds of hours of unindexed lecture video every
day. Manual chaptering does not scale, and existing automatic methods are
either text-only, closed-source, or restricted to a single (flat) level of
segmentation. This thesis presents \lecseg{}, the first open hierarchical
multimodal pipeline for lecture-video topic segmentation that produces both
chapter-level and subtopic-level boundaries.

We make four technical contributions. (1) A reliability-weighted modality-fusion
module learns to down-weight unreliable channels (e.g.\ noisy OCR on a chalkboard
lecture) per video and per sentence. (2) An explicit two-level hierarchy is
modelled jointly, rather than as two independent flat tasks. (3) Boundary
refinement and chapter titling are performed by a small, locally-hosted open
language model, removing the dependence on closed APIs reported in prior work.
(4) We release \dataset{}: a corpus of 30 lecture videos covering 5 academic
domains, with dual subtopic annotation and reported inter-annotator agreement.

We evaluate against three classical and two neural baselines on five metrics:
$P_k$, WindowDiff, Boundary Similarity, tolerance-F\textsubscript{1}, and a
hierarchical WindowDiff variant we propose. All results are reported with
1{,}000-sample bootstrap confidence intervals and pairwise Wilcoxon
signed-rank tests at video granularity. \lecseg{} reduces $P_k$ by
\todo{X.XX}\,\% relative ($\Delta=$\,\todo{0.0XX} absolute) over the strongest
baseline, with $p < \todo{0.0X}$. The full reproduction artefact regenerates
every numerical claim of this thesis from a fresh clone with a single command
(\texttt{make reproduce}).

\vspace{1em}
\noindent\textbf{Keywords:} topic segmentation, lecture video, multimodal
fusion, hierarchical segmentation, local language models, reproducible
research.
